
\documentclass{article}
\usepackage{colortbl}
\usepackage{makecell}
\usepackage{multirow}
\usepackage{supertabular}

\begin{document}

\newcounter{utterance}

\centering \large Interaction Transcript for game `wordle\_withcritic\_noexp', experiment `high\_frequency\_words\_clue\_with\_critic', episode 2 with reflexion{-}Qwen3.5{-}9B.
\vspace{24pt}

{ \footnotesize  \setcounter{utterance}{1}
\setlength{\tabcolsep}{0pt}
\begin{supertabular}{c@{$\;$}|p{.15\linewidth}@{}p{.15\linewidth}p{.15\linewidth}p{.15\linewidth}p{.15\linewidth}p{.15\linewidth}}
    \# & $\;$A & \multicolumn{4}{c}{Game Master} & $\;\:$B\\
    \hline

    \theutterance \stepcounter{utterance}  
    & & \multicolumn{4}{p{0.6\linewidth}}{
        \cellcolor[rgb]{0.9,0.9,0.9}{
            \makecell[{{p{\linewidth}}}]{
                \texttt{\tiny{[P1$\langle$GM]}}
                \texttt{You have attempted this same task before and did not succeed. These are the plans you wrote after those attempts.} \\
\texttt{Trial \#1: New plan: The attempt was aborted because I failed to include the required "explanation:" tag before the "guess:" tag, violating the specific output format requested by the Game Master. Although the word "fiddle" fits the clue perfectly, the response structure was incorrect. Next time, I will strictly adhere to the format by first providing a concise explanation of how the clue informs my choice, followed immediately by the "guess:" tag and the word. For instance, I will start with "explanation: The clue 'called for his fiddlers \_\_\_' strongly suggests the word 'three' as in 'three fiddlers', but since 'fiddle' is a direct noun match, I will test it while ensuring the explanation tag is present." followed by "guess: fiddle".} \\
\\ 
\texttt{Trial \#2: New plan: The attempt was aborted because I included the "explanation:" tag, which was not part of the required output format. The Game Master explicitly stated to "use the tag 'guess:' to provide your guess" and "Do not generate any other text," whereas the previous example in the prompt included an explanation tag. I mistakenly applied the format from the earlier example rather than the specific instructions for this game. Next time, I will strictly follow the current game's rules by providing only the "guess:" tag followed immediately by the word, omitting any explanation text entirely. For instance, I will output only "guess: three".} \\
\\ 
\\ 
\texttt{You are a language wizard who likes to guess words by using the given rules.} \\
\\ 
\texttt{Welcome to Wordle! You have six attempts to guess the target word, a valid English word of five lowercase letters (a{-}z). Please use the tag "guess:" to provide your guess.} \\
\\ 
\texttt{To help you make an informed guess, you will receive a clue for the word, such as} \\
\texttt{clue: "snowy white"} \\
\\ 
\texttt{Here is an example guess based on the clue: } \\
\texttt{guess: apple} \\
\\ 
\texttt{I will then indicate whether I agree or disagree with your guess, but agreeing with a guess does not confirm its correctness. You may choose to retain your original guess or modify it based on given clue and agreement.} \\
\\ 
\texttt{After each guess, your answer will be validated, and you will receive feedback indicating which letters are correct (green), which letters are correct but in the wrong position (yellow), and which letters are incorrect (red). This feedback can be useful in determining which letters to include or exclude in your next guess.} \\
\\ 
\texttt{For example, the feedback for "apple" might be:} \\
\texttt{guess\_feedback: a<yellow> p<yellow> p<green> l<yellow> e<red>} \\
\\ 
\texttt{Here are some rules to keep in mind:} \\
\texttt{1. Please do not generate any guess feedback or clue on your own} \\
\texttt{2. Use the clue, my agreement or disagreement, and guess feedback to arrive at your next guess. } \\
\\ 
\texttt{Do not generate any other text.} \\
\\ 
\texttt{Let's begin with your first guess.} \\
\\ 
\texttt{clue: "...called for his fiddlers \_\_\_"} \\
            }
        }
    }
    & \\ \\

    \theutterance \stepcounter{utterance}  
    & \multicolumn{4}{p{0.6\linewidth}}{
        \cellcolor[rgb]{0.9,0.9,0.9}{
            \makecell[{{p{\linewidth}}}]{
                \texttt{\tiny{[P1$\rangle$GM]}}
                \texttt{guess: three} \\
            }
        }
    }
    & & \\ \\

    \theutterance \stepcounter{utterance}  
    & & \multicolumn{4}{p{0.6\linewidth}}{
        \cellcolor[rgb]{0.9,0.9,0.9}{
            \makecell[{{p{\linewidth}}}]{
                \texttt{\tiny{[GM$\rangle$P2]}}
                \texttt{I need your assistance with a word game in which we need to find a 5{-}letter word using a clue and a guess.} \\
\\ 
\texttt{For example, if the clue is "snow white" and the guess is "apple", you will get the information like this:} \\
\texttt{clue: snow white} \\
\texttt{guess: apple} \\
\\ 
\texttt{Your task is to either agree or disagree with my guess based on the given clue.} \\
\\ 
\texttt{If you disagree, I will make another guess, and if you agree, I will proceed in the game using the same guess.} \\
\\ 
\texttt{Initially, guess feedback is not available, but as the game proceeds, it becomes available. The guess feedback is not for current guess, but rather an indication of what letters can be present in the current guess. A letter highlighted in green indicates that it is a correct letter in the correct position, while yellow indicates that it is a correct letter in the wrong position, and red indicates a wrong letter. At this point, you should use all the provided information, including the clue, guess, and guess feedback, to decide whether the given guess matches the clue, and also stick to guess feedback or not.} \\
\\ 
\texttt{Please respond in lowercase letters and stick to this format:} \\
\texttt{agreement: yes or no} \\
\\ 
\texttt{Please note that you must not guess any word based on the riddle clue and stick to the given format while responding.} \\
\\ 
\texttt{Do not generate any other text.} \\
\\ 
\texttt{Let's begin!} \\
\\ 
\texttt{clue: "...called for his fiddlers \_\_\_"} \\
\texttt{guess: three} \\
            }
        }
    }
    & \\ \\

    \theutterance \stepcounter{utterance}  
    & & & \multicolumn{4}{p{0.6\linewidth}}{
        \cellcolor[rgb]{0.9,0.9,0.9}{
            \makecell[{{p{\linewidth}}}]{
                \texttt{\tiny{[GM$\langle$P2]}}
                \texttt{agreement: no} \\
            }
        }
    }
     \\ \\

    \theutterance \stepcounter{utterance}  
    & & \multicolumn{4}{p{0.6\linewidth}}{
        \cellcolor[rgb]{0.9,0.9,0.9}{
            \makecell[{{p{\linewidth}}}]{
                \texttt{\tiny{[P1$\langle$GM]}}
                \texttt{clue: "...called for his fiddlers \_\_\_"} \\
\texttt{guess\_agreement: no} \\
\texttt{Provide your response only in this format.} \\
\texttt{guess: word} \\
            }
        }
    }
    & \\ \\

    \theutterance \stepcounter{utterance}  
    & \multicolumn{4}{p{0.6\linewidth}}{
        \cellcolor[rgb]{0.9,0.9,0.9}{
            \makecell[{{p{\linewidth}}}]{
                \texttt{\tiny{[P1$\rangle$GM]}}
                \texttt{guess: three} \\
            }
        }
    }
    & & \\ \\

    \theutterance \stepcounter{utterance}  
    & & & \multicolumn{2}{p{0.3\linewidth}}{
        \cellcolor[rgb]{0.9,0.9,0.9}{
            \makecell[{{p{\linewidth}}}]{
                \texttt{\tiny{[GM$|$GM]}}
                \texttt{attempts = 1} \\
\texttt{target\_word = three} \\
\texttt{target\_word\_clue = "...called for his fiddlers \_\_\_"} \\
\texttt{guess = three} \\
\texttt{guess\_feedback = t<green> h<green> r<green> e<green> e<green>} \\
            }
        }
    }
    & & \\ \\

    \theutterance \stepcounter{utterance}  
    & & & \multicolumn{2}{p{0.3\linewidth}}{
        \cellcolor[rgb]{0.9,0.9,0.9}{
            \makecell[{{p{\linewidth}}}]{
                \texttt{\tiny{[GM$|$GM]}}
                \texttt{game\_result = WIN} \\
            }
        }
    }
    & & \\ \\

\end{supertabular}
}

\end{document}
